\documentclass[11pt]{letter}
\usepackage[margin=1in]{geometry}
\usepackage[T1]{fontenc}
\usepackage{parskip}
\signature{Ku Kang, on behalf of all authors}
\address{Ku Kang\\CBRN Defense Research Institute\\Republic of Korea\\bisu9082@gmail.com}
\begin{document}
\begin{letter}{The Editors\\Journal of Hazardous Materials}
\opening{Dear Editors,}

We submit the research article \textit{``A testing-effort artefact in the
most-sensitive aquatic endpoint: structure-based PFAS screening recovers
predictive validity but still cannot support clearance''} for consideration in
\textit{Journal of Hazardous Materials}.

\textbf{Is the studied subject an environmental contaminant?} The work concerns
per- and polyfluoroalkyl substances, analysed through 15{,}991 measured
freshwater aquatic toxicity records for 103 substances drawn from the U.S. EPA
ECOTOX Knowledgebase, and applied to 17{,}619 untested structures from the EPA
PFASSTRUCTV6 inventory. The broader diagnostic uses 106{,}095 freshwater acute
records across 4{,}757 chemicals.

\textbf{Is the study conducted under environmentally relevant conditions?} The
modelled quantity is a measured effect concentration in freshwater in mass
units, and the decision analysis is evaluated against citable hazard
boundaries rather than an arbitrary cut-off: we report clearance across a
continuum of thresholds and at the Globally Harmonised System acute aquatic
Category 1 boundary of 1 mg/L, as well as at a threshold matched for stringency
to the protective convention.

\textbf{The contribution.} Regulatory hazard screening conventionally summarises
a substance by the lowest effect concentration reported for it. We show that
this makes the response variable depend on how often the substance has been
tested: across PFAS the protective minimum correlates with record count at
$\rho=-0.73$, and subsampling to equal depth identifies about half of that as a
pure extreme-order-statistic artefact. Because a structure-based model cannot
predict testing history, this confound is by itself sufficient to produce the
external validation failures reported in the PFAS screening literature. Holding
a published model fixed and changing only how the same records are summarised
moves external performance from $R^2=-0.75$ to $+0.22$. Rebuilding training and
external sets consistently on an effort-stable endpoint restores external
validity and halves the prediction interval --- and yet no untested structure
clears at a stringency-matched threshold.

We believe this is useful to the journal's readership in two ways. It identifies
a measurable defect in a widely used data-preparation convention, with a
diagnostic that any screening study can compute in one line. And it corrects the
interpretation of a negative result that is becoming established: structure-based
models are not failing on PFAS because they cannot represent fluorinated
chemistry, but substantially because they are being asked to predict a quantity
that encodes testing effort.

\textbf{Prior submission.} An earlier and materially different version of this
work, based on a curated subset of the aquatic record and without the
testing-effort analysis, was reviewed and declined by \textit{Environmental
Science \& Technology} (es-2026-093898). The present manuscript uses a different
data source, reports a different central finding, and reaches its conclusion by
a different mechanism; we mention the prior review in the interest of
transparency and can supply those reviews on request.

All extraction and analysis scripts, structure-locked tables and inventory-level
predictions are openly available and tagged as an immutable release. The
manuscript is original, is not under consideration elsewhere, and is approved by
all authors. The authors declare no competing financial interest.

\closing{Sincerely,}
\end{letter}
\end{document}
